\documentclass[11pt]{article}
\usepackage[T1]{fontenc}
\usepackage[utf8]{inputenc}
\usepackage[spanish]{babel}
\usepackage{amsmath,amssymb}
\usepackage{booktabs}
\usepackage{array}
\usepackage{adjustbox}
\usepackage{longtable}
\usepackage{xcolor}
\usepackage{fancyvrb}
\usepackage[margin=2.2cm]{geometry}
\usepackage{hyperref}

\newcommand{\col}[1]{\texttt{#1}}
\newcommand{\archivo}[1]{\textbf{\texttt{#1}}}
\fvset{fontsize=\footnotesize,frame=single,framesep=2mm}

\title{Bases de datos del modelo y del Monte Carlo\\
\large Qué se genera, cómo leerlo y qué gráficas salen}
\author{Notas de Claude para Rodrigo}
\date{18 de noviembre de 2026}

\sloppy
\begin{document}
\maketitle
\tableofcontents

\section{Panorama}

Todo se escribe en \texttt{claude/niu/final/output/modelo\_tesis/}, en dos carpetas:

\begin{center}
\begin{adjustbox}{max width=\textwidth}\begin{tabular}{lll}
\toprule
Carpeta & La genera & Qué hay\\
\midrule
\texttt{teoria/<tag>/} & \texttt{teoria.py} & El equilibrio \emph{verdadero} de los 3 regímenes y un panel simulado.\\
\texttt{montecarlo/<tag>/} & \texttt{montecarlo.py}, \texttt{por\_edad\_mc.py} & Las estimaciones de cada réplica y el equilibrio en cada $\hat\theta$.\\
\bottomrule
\end{tabular}\end{adjustbox}
\end{center}

\paragraph{Orden para correr en GPU.}
\begin{Verbatim}
cd claude/niu/final/modelo_tesis
python -u teoria.py --tag R_LB --guardar_panel -v       # 1. equilibrio verdadero
python -u montecarlo.py --reps 0:50                     # 2. estimaciones (o en dos GPUs)
python -u montecarlo.py --summarize                     # 3. tabla resumen
python -u por_edad_mc.py --reps 0:50                    # 4. equilibrio en cada theta estimado
\end{Verbatim}
\texttt{montecarlo.py} y \texttt{por\_edad\_mc.py} buscan el equilibrio verdadero en
\texttt{teoria/R\_LB/}. Por eso el paso 1 usa \texttt{--tag R\_LB}; si usas otro tag, pásalo
con \texttt{--desde <tag>}.

\paragraph{Convenciones que se repiten en todos los archivos.}
\begin{itemize}
  \item \col{regimen}: 0, 1, 2 $\leftrightarrow$ $\zeta=-0.3,\,0,\,+0.3$ (reparación barata,
  base, cara). El régimen 1 es el base.
  \item \col{j}: marca. 0 = light brown, 1 = heavy brown.
  \item \col{a}: edad del coche, $1,\dots,24$. En el panel, $a=0$ es terminal y $a=-1$ sin coche.
  \item \col{tipo}, o \col{t} en los nombres de parámetros: 0 = pareja pobre, 1 = soltero pobre.
  \item Precios en miles de coronas danesas (DKK). Utilidades en utils.
  \item \col{q}: fracción de la masa de hogares. Suma 1 sobre \emph{todos} los estados,
  incluidos sin coche y terminal, así que $\sum_{j,a}q<1$.
  \item Todo se promedia sobre los dos tipos con pesos $f=(1/2,1/2)$, salvo que se diga otra cosa.
\end{itemize}

\paragraph{Sobre los ejemplos.} Las filas de \texttt{teoria/} vienen de una corrida
\emph{de prueba} con la especificación nueva (\texttt{teoria.py --prueba}: $\bar a=10$, grid
de $w$ chico). Las del Monte Carlo vienen del MC de octubre (\texttt{claude/niu/output/}), que
tiene el mismo formato. Con la especificación nueva solo cambian dos cosas: ya no existe el
parámetro \col{tc\_sell\_inspect}, y \col{delta} ahora vale 0.075. Los números de los
ejemplos son ilustrativos y están redondeados; los archivos traen precisión completa.

\paragraph{Leer en R.}
\begin{Verbatim}
library(tidyverse)
base <- "claude/niu/final/output/modelo_tesis"
mc   <- file.path(base, "montecarlo", "base")
par  <- list.files(mc, "^parametros_reps.*\\.csv$", full.names = TRUE) |>
          map_dfr(read_csv) |> filter(mejor)
\end{Verbatim}
En Python: \texttt{pd.concat(pd.read\_csv(f) for f in glob.glob(...))}.

% ===========================================================================
\section{Equilibrio verdadero: \texttt{teoria/<tag>/}}

\subsection{\archivo{resumen.csv}: agregados de mercado por régimen}

Hay una fila por régimen y una más con Gillingham resuelto en los mismos parámetros
(\col{regimen = gillingham}).

\begin{center}\small
\begin{adjustbox}{max width=\textwidth}\begin{tabular}{ll}
\toprule
Columna & Significado\\
\midrule
\col{zeta} & $\zeta$ del régimen ($R_m=R\,e^{\zeta}$)\\
\col{ok}, \col{iters}, \col{segundos} & convergencia y tiempo del solver\\
\col{sin\_coche} & fracción de hogares sin coche\\
\col{compra\_nuevo}, \col{compra\_usado} & fracción de hogares que compran un nuevo / un usado al año\\
\col{reparacion} & reparaciones por coche en circulación al año\\
\col{chatarreo\_endogeno} & chatarreos voluntarios por coche al año\\
\col{accidentes} & accidentes por coche al año\\
\col{edad\_media} & edad media del parque\\
\col{w\_medio} & desgaste medio del parque\\
\col{masa\_borde\_w} & masa en los dos extremos del grid de $w$ (debe ser $<0.5\%$)\\
\bottomrule
\end{tabular}\end{adjustbox}
\end{center}

\begin{center}\footnotesize
\begin{adjustbox}{max width=\textwidth}\begin{tabular}{rrrrrrrrrr}
\toprule
regimen & zeta & sin\_coche & compra\_nuevo & compra\_usado & reparacion & chatarreo\_end. & accidentes & w\_medio & masa\_borde\_w\\
\midrule
0 & $-0.3$ & 0.435 & 0.059 & 0.187 & 0.182 & 0.0072 & 0.0094 & 0.136 & 0.0044\\
1 & $0.0$ & 0.442 & 0.058 & 0.188 & 0.109 & 0.0072 & 0.0103 & 0.232 & 0.0057\\
2 & $0.3$ & 0.447 & 0.058 & 0.188 & 0.052 & 0.0071 & 0.0110 & 0.303 & 0.0076\\
\bottomrule
\end{tabular}\end{adjustbox}
\end{center}
Se lee así: cuando la reparación es más cara, se repara menos, el parque está más desgastado
y hay más accidentes. Este es el mecanismo de la variable excluida. (En la prueba el grid es
chico y \col{masa\_borde\_w} pasa de 0.5\%; en la corrida completa hay que revisarlo.)

\subsection{\archivo{por\_edad.csv}: el equilibrio por marca y edad}

Una fila por régimen $\times$ marca $\times$ edad. Ya está promediado sobre $w$ y tipos.

\begin{center}\small
\begin{adjustbox}{max width=\textwidth}\begin{tabular}{p{3.3cm}p{11.5cm}}
\toprule
Columna & Significado\\
\midrule
\col{q} & masa de hogares que empieza el año con un coche $(j,a)$\\
\col{P\_stock} & precio medio, ponderado por los coches que existen (sobre $w$)\\
\col{P\_transaccion} & precio medio, ponderado por los coches que se ofrecen en venta (el que ``vería'' un comprador)\\
\col{keep} & Pr(mantener)\\
\col{reparar} & Pr(reparar) promedio (0 en la edad $\bar a-1$)\\
\col{w\_medio} & desgaste medio de los coches de esa edad\\
\col{q\_gill}, \col{P\_gill}, \col{keep\_gill} & lo mismo en el modelo de Gillingham con los mismos parámetros (igual en los 3 regímenes)\\
\bottomrule
\end{tabular}\end{adjustbox}
\end{center}

\begin{center}\footnotesize
\begin{adjustbox}{max width=\textwidth}\begin{tabular}{rrrrrrrrrrrr}
\toprule
regimen & j & a & q & P\_stock & P\_trans. & keep & reparar & w\_medio & q\_gill & P\_gill & keep\_gill\\
\midrule
0 & 0 & 1 & 0.0512 & 137.58 & 137.63 & 0.685 & 0.225 & 0.075 & 0.0514 & 136.73 & 0.704\\
0 & 0 & 2 & 0.0510 & 121.45 & 121.52 & 0.678 & 0.253 & 0.093 & 0.0512 & 119.99 & 0.703\\
0 & 0 & 3 & 0.0507 & 105.52 & 105.58 & 0.674 & 0.267 & 0.105 & 0.0509 & 103.67 & 0.702\\
\bottomrule
\end{tabular}\end{adjustbox}
\end{center}
En la edad 1, \col{w\_medio}$=0.075=\delta$: el coche nació en $w=0$ y le tocó un año de
desgaste sin reparación (los nuevos no se reparan).

\subsection{\archivo{panel\_resumen.csv}: lo que se ve en el panel simulado}

\begin{center}\footnotesize
\begin{adjustbox}{max width=\textwidth}\begin{tabular}{rrrrrrr}
\toprule
regimen & hogar\_anios & sin\_coche & tasa\_reparacion & accidentes\_por\_coche & chatarreo\_endogeno & compra\\
\midrule
0 & 8000 & 0.437 & 0.223 & 0.0102 & 0.0073 & 0.250\\
1 & 8000 & 0.449 & 0.140 & 0.0116 & 0.0073 & 0.248\\
2 & 8000 & 0.457 & 0.062 & 0.0120 & 0.0067 & 0.245\\
\bottomrule
\end{tabular}\end{adjustbox}
\end{center}
Sirve para validar que el panel reproduce al modelo. Ojo: aquí \col{tasa\_reparacion} es la
tasa entre los coches que \emph{pueden} repararse; en \texttt{resumen.csv} es por coche en
circulación. Por eso los números no coinciden.

\subsection{\archivo{panel.csv.gz}: los microdatos (solo con \texttt{--guardar\_panel})}

Una fila por hogar y año. Es lo único que tiene la reparación verdadera de cada coche.

\begin{center}\small
\begin{longtable}{p{3.6cm}p{11.2cm}}
\toprule
Columna & Significado\\
\midrule
\col{regimen}, \col{tipo}, \col{id}, \col{k} & hogar (\col{id} es único dentro de cada régimen) y año $k=0,\dots,3$\\
\col{j}, \col{a}, \col{w} & estado al \emph{inicio} del año, antes de comerciar. $a=0$: terminal; $j=a=w=-1$: sin coche\\
\col{o} & qué pasó con el estado inicial (ver abajo)\\
\col{j\_h}, \col{d\_h}, \col{w\_h} & tenencia después de comerciar: el coche que maneja este año. $d_h=0$: nuevo; $-1$: sin coche\\
\col{r} & reparó la tenencia este año: 1/0; $-1$ si no aplica (sin coche, nuevo o edad $\bar a-1$)\\
\col{accidente} & la tenencia se descompuso este año\\
\col{chatarreo\_endogeno} & se deshizo de un coche activo y lo chatarrizó por decisión propia\\
\col{x}, \col{h} & los mismos estados como índice en el vector de estados (solo para la verosimilitud)\\
\bottomrule
\end{longtable}
\end{center}

\paragraph{Ojo: \col{w} y \col{w\_h} son el \emph{índice} del grid, no el valor.} El valor es
\[
w = \texttt{w\_min} + \texttt{índice}\times\texttt{w\_step},
\]
que en la corrida completa es $-3+0.075\times\text{índice}$ (índice $0,\dots,70$; $w=0$ es el
índice 40). En la prueba es $-1.5+0.15\times\text{índice}$. Los valores están en
\texttt{config.json}.

\paragraph{Códigos de \col{o}.}
\begin{center}\small
\begin{adjustbox}{max width=\textwidth}\begin{tabular}{clll}
\toprule
\col{o} & Coche activo & Coche terminal & Sin coche\\
\midrule
0 & mantener & --- & ---\\
1 & vender y quedarse sin coche & --- & quedarse sin coche\\
2 & chatarrizar y quedarse sin coche & chatarra y sin coche & ---\\
3 & vender y comprar otro & --- & comprar\\
4 & chatarrizar y comprar otro & chatarra y comprar & ---\\
\bottomrule
\end{tabular}\end{adjustbox}
\end{center}

\begin{center}\footnotesize
\begin{adjustbox}{max width=\textwidth}\begin{tabular}{rrrrrrrrrrrrr}
\toprule
regimen & tipo & id & k & j & a & w & o & j\_h & d\_h & w\_h & r & accidente\\
\midrule
0 & 0 & 0 & 0 & $-1$ & $-1$ & $-1$ & 1 & $-1$ & $-1$ & $-1$ & $-1$ & 0\\
0 & 0 & 0 & 1 & $-1$ & $-1$ & $-1$ & 3 & 1 & 1 & 11 & 0 & 0\\
0 & 0 & 0 & 2 & 1 & 2 & 13 & 1 & $-1$ & $-1$ & $-1$ & $-1$ & 0\\
0 & 0 & 0 & 3 & $-1$ & $-1$ & $-1$ & 3 & 0 & 3 & 15 & 0 & 0\\
\bottomrule
\end{tabular}\end{adjustbox}
\end{center}
La historia del hogar 0:
\begin{itemize}
  \item Año 0: no tiene coche y sigue sin coche.
  \item Año 1: compra un heavy brown de 1 año ($w$ índice 11) y no lo repara.
  \item Año 2: empieza con ese coche, ahora de 2 años con $w$ índice 13 (el desgaste subió),
  lo vende y se queda sin coche.
  \item Año 3: compra un light brown de 3 años.
\end{itemize}
La transición de un año a otro es: la tenencia $(j_h,d_h,w_h)$ y la decisión $r$ del año $k$
dan el estado $(j,a,w)$ del año $k+1$, salvo que haya habido accidente.

\subsection{Otros archivos de \texttt{teoria/}}
\begin{itemize}
  \item \archivo{gill\_por\_edad.csv} (\col{j, a, q, P, keep}) y
  \archivo{gill\_resumen.json}: Gillingham en los mismos parámetros. Ya vienen pegados en
  \texttt{por\_edad.csv} y \texttt{resumen.csv}.
  \item \archivo{equilibrios.npz}, \archivo{config.json}: precios y $EV$ completos de cada
  régimen, y parámetros. No son para graficar directamente. Si quieres el precio por $w$
  (no solo el promedio por edad), pídemelo y lo exporto a CSV.
\end{itemize}

% ===========================================================================
\section{Monte Carlo: \texttt{montecarlo/<tag>/}}

Los archivos llevan el rango de réplicas en el nombre (\texttt{reps0-24}, \texttt{reps25-49},
\dots) porque se corren en bloques, en dos GPUs. Hay que juntarlos. Los que terminan en
\texttt{\_faltantes} son re-estimaciones de las que fallaron: si un par (réplica, estimador)
aparece dos veces, quédate con el del archivo \texttt{\_faltantes}.

\paragraph{Estimadores} (columna \col{estimador}):
\begin{center}\small
\begin{adjustbox}{max width=\textwidth}\begin{tabular}{lll}
\toprule
\col{estimador} & Modelo & Qué observa\\
\midrule
\col{diseno\_1} & tesis & $r$ y el motivo de salida (cota superior)\\
\col{diseno\_2} & tesis & el motivo de salida, no $r$\\
\col{diseno\_3} & tesis & ni $r$ ni el motivo (como los datos reales)\\
\col{gill\_parcial} & Gillingham & no ve $w$, $r$ ni el motivo (la verosimilitud del paper)\\
\col{gill\_completa} & Gillingham & no ve $w$ ni $r$; sí el motivo\\
\bottomrule
\end{tabular}\end{adjustbox}
\end{center}

\subsection{\archivo{parametros\_reps*.csv} y \archivo{gill\_parametros\_reps*.csv}: la base principal}

Una fila por réplica $\times$ estimador $\times$ arranque $\times$ parámetro.

\begin{center}\small
\begin{adjustbox}{max width=\textwidth}\begin{tabular}{ll}
\toprule
Columna & Significado\\
\midrule
\col{rep} & réplica (también es la semilla)\\
\col{arranque}, \col{mejor} & punto inicial; \col{mejor = True} es la estimación que cuenta. \textbf{Filtra siempre por \col{mejor}.}\\
\col{parametro} & nombre (ver abajo)\\
\col{verdad}, \col{estimado}, \col{se} & valor verdadero, estimación y error estándar (BHHH), en escala natural\\
\col{ll}, \col{convergio} & log-verosimilitud y convergencia (iguales en todas las filas de una estimación)\\
\bottomrule
\end{tabular}\end{adjustbox}
\end{center}

\begin{center}\footnotesize
\begin{adjustbox}{max width=\textwidth}\begin{tabular}{rllrlrrrrl}
\toprule
rep & estimador & arranque & mejor & parametro & verdad & estimado & se & ll & convergio\\
\midrule
1 & diseno\_1 & 0 & True & mu\_t0 & 0.1131 & 0.1110 & 0.0030 & $-357598.6$ & True\\
1 & diseno\_1 & 0 & True & mu\_t1 & 0.0941 & 0.0929 & 0.0027 & $-357598.6$ & True\\
1 & diseno\_1 & 0 & True & u0\_t0\_j0 & 3.6490 & 3.6255 & 0.0545 & $-357598.6$ & True\\
1 & gill\_parcial & 0 & True & mu\_t0 & 0.1131 & 0.0511 & 0.0049 & $-145797.9$ & True\\
\bottomrule
\end{tabular}\end{adjustbox}
\end{center}
La última fila muestra el sesgo de Gillingham: estima $\mu$ en menos de la mitad del valor
verdadero.

\paragraph{Los 26 parámetros del modelo de la tesis} (los nombres siguen el patrón
\col{nombre\_t<tipo>\_j<marca>}):
\begin{center}\small
\begin{adjustbox}{max width=\textwidth}\begin{tabular}{lll}
\toprule
Grupo & Nombres en el CSV & Tesis\\
\midrule
Preferencias & \col{mu\_t0, mu\_t1} & $\mu_\tau$\\
 & \col{u0\_t0\_j0, u0\_t0\_j1, u0\_t1\_j0, u0\_t1\_j1} & $u_{0,\tau j}$\\
 & \col{u1\_t0\_j0, \dots, u1\_t1\_j1} & $u_{1,\tau j}$\\
 & \col{u\_w} & $u_w$\\
Transacción & \col{tc\_buy\_t0, tc\_buy\_t1} & $T_{b,\tau}$\\
 & \col{tc\_buy\_nocar\_t0, tc\_buy\_nocar\_t1} & $T^\emptyset_{b,\tau}$\\
 & \col{tc\_sell}, \col{sigma\_sell} & $T_s$, $\sigma_D$\\
Accidentes & \col{acc\_int\_j0, acc\_int\_j1}, \col{acc\_age\_j0, acc\_age\_j1} & $\alpha_{0j}$, $\alpha_{1j}$\\
Reparación & \col{delta\_j0, delta\_j1}, \col{kappa}, \col{sigma\_eta}, \col{sigma\_rep} & $\delta_j$, $\kappa$, $\sigma_\eta$, $\sigma_r$\\
\bottomrule
\end{tabular}\end{adjustbox}
\end{center}
Gillingham comparte 20 (todos menos \col{u\_w} y los 5 de reparación). Esos 20 son
los que se comparan entre los cinco estimadores.

\subsection{\archivo{resumen\_reps*.csv} y \archivo{gill\_resumen\_reps*.csv}: una fila por estimación}

\begin{center}\small
\begin{adjustbox}{max width=\textwidth}\begin{tabular}{p{4.2cm}p{10.6cm}}
\toprule
Columna & Significado\\
\midrule
\col{ll}, \col{n\_obs}, \col{n\_celdas} & log-verosimilitud, transiciones (90,000) y celdas distintas\\
\col{convergio}, \col{n\_eval}, \col{n\_fail}, \col{segundos} & diagnóstico del optimizador\\
\col{cuerda}, \col{refact}, \col{newton}, \col{cond\_B} & diagnóstico del solver de equilibrio y número de condición de BHHH (no son para la tesis)\\
\col{P\_rmse} & error cuadrático medio de los precios de equilibrio en $\hat\theta$ contra los verdaderos (miles de DKK, ponderado por la masa de coches)\\
\col{reparacion\_verdad}, \col{reparacion\_est} & tasa agregada de reparación verdadera y la que implica $\hat\theta$\\
\midrule
\multicolumn{2}{l}{\textit{Solo en \texttt{gill\_resumen}:}}\\
\col{sin\_coche\_est}, \col{compra\_nuevo\_est}, \dots & los agregados de mercado que implica el $\hat\theta$ de Gillingham (se comparan con \texttt{teoria/resumen.csv})\\
\bottomrule
\end{tabular}\end{adjustbox}
\end{center}

\begin{center}\footnotesize
\begin{adjustbox}{max width=\textwidth}\begin{tabular}{rlrrrrrr}
\toprule
rep & estimador & ll & n\_celdas & segundos & P\_rmse & reparacion\_verdad & reparacion\_est\\
\midrule
1 & diseno\_1 & $-357598.6$ & 48713 & 139.6 & 0.464 & 0.2760 & 0.2770\\
1 & diseno\_2 & $-338913.2$ & 45753 & 124.2 & 0.529 & 0.2760 & 0.2751\\
1 & diseno\_3 & $-334719.5$ & 45239 & 99.7 & 0.744 & 0.2760 & 0.2756\\
\bottomrule
\end{tabular}\end{adjustbox}
\end{center}
Las \col{ll} no se comparan entre diseños, porque cada diseño tiene datos distintos.

\subsection{\archivo{resumen\_mc.csv}: la tabla ya resumida (\texttt{--summarize})}

Una fila por estimador $\times$ parámetro, con columnas \col{verdad, media, sesgo, sesgo\_rel,
sd, se\_mediano, rmse, cobertura, reps}. La \col{cobertura} es la fracción de réplicas en que
$\hat\theta\pm1.96\,se$ contiene la verdad. Ejemplo (de una prueba de una sola réplica, por
eso \col{sd} está vacía):
\begin{center}\footnotesize
\begin{adjustbox}{max width=\textwidth}\begin{tabular}{llrrrrrrrr}
\toprule
estimador & parametro & verdad & media & sesgo & sesgo\_rel & se\_mediano & rmse & cobertura & reps\\
\midrule
diseno\_1 & mu\_t0 & 0.1131 & 0.0953 & $-0.0178$ & $-0.157$ & 0.0158 & 0.0178 & 1.0 & 1\\
\bottomrule
\end{tabular}\end{adjustbox}
\end{center}
Si haces tus propias tablas desde \texttt{parametros\_reps*.csv}, no necesitas este archivo.

\subsection{\archivo{por\_edad\_mc.csv} y \archivo{gill\_por\_edad\_mc.csv}: el equilibrio en cada $\hat\theta$}

Son las mismas columnas que \texttt{teoria/por\_edad.csv} (sin las de Gillingham), más
\col{rep} y \col{estimador}. Las filas con \col{rep = -1} y \col{estimador = verdad} (o
\col{gill\_verdad}) son la verdad. La versión de Gillingham no tiene régimen ni reparación:
\col{rep, estimador, j, a, q, P, keep}.

\begin{center}\footnotesize
\begin{adjustbox}{max width=\textwidth}\begin{tabular}{rlrrrrrrrrr}
\toprule
rep & estimador & regimen & j & a & q & P\_stock & P\_trans. & keep & reparar & w\_medio\\
\midrule
$-1$ & verdad & 0 & 0 & 1 & 0.0291 & 140.49 & 140.53 & 0.672 & 0.443 & 0.075\\
$-1$ & verdad & 0 & 0 & 2 & 0.0289 & 128.70 & 128.77 & 0.663 & 0.498 & 0.039\\
\bottomrule
\end{tabular}\end{adjustbox}
\end{center}
Con esto se dibuja, para cada edad, la verdad contra la banda de las 50 réplicas: precios,
reparación y distribución \emph{implicados} por cada estimador.

\subsection{\archivo{fallas\_reps*.csv}}
\col{rep, diseno, arranque, error}: las estimaciones que no terminaron. Sirve para reportar
cuántas réplicas hay por estimador.

\subsection{\archivo{inferencia\_reparacion\_reps*.csv.gz}: qué tanto se infiere $r$ de cada coche}

Este archivo responde la pregunta de la tesis observación por observación. Cada observación
es un coche que un hogar tuvo en el año $k$ y que se \emph{podía} reparar (usado de edad
$1,\dots,\bar a-2$), junto con lo que pasó en el año $k+1$. Para cada una, el archivo trae la
reparación verdadera y la probabilidad de que haya reparado según el modelo, con tres
niveles de información:

\begin{center}\small
\begin{adjustbox}{max width=\textwidth}\begin{tabular}{lp{6.2cm}p{6.2cm}}
\toprule
Columna & Qué es & Fórmula\\
\midrule
\col{prior} & Pr(reparar) antes de ver el año siguiente: la CCP & $p=\Pi(r=1\mid j,a,w)$\\
\col{post\_d2} & después de ver lo que ve el diseño 2: el $w'$ del coche si sobrevive & $\dfrac{p\,f(w'\mid r=1)}{p\,f(w'\mid r=1)+(1-p)\,f(w'\mid r=0)}$; si hubo accidente, $=p$\\
\col{post\_d3} & después de ver lo que ve el diseño 3: si el coche se chatarrizó no se ve ni el $w'$ ni si fue accidente & la misma regla de Bayes con la mezcla de la verosimilitud del diseño 3; si no se chatarrizó, $=$ \col{post\_d2}\\
\bottomrule
\end{tabular}\end{adjustbox}
\end{center}

Es la regla de Bayes aplicada a la mezcla $Z$ de la verosimilitud. Las probabilidades se
calculan con cada estimador (\col{estimador}): \col{verdad} usa $\theta$ verdadero (lo mejor
que se puede inferir) y \col{diseno\_1/2/3} usan el $\hat\theta$ de ese diseño.
Para estudiar el diseño 2, usa \col{estimador = diseno\_2} con \col{post\_d2}; para el
diseño 3, \col{diseno\_3} con \col{post\_d3}.

\paragraph{Está agregado por celdas.} Todas las observaciones con el mismo régimen, tipo,
coche $(j,a,w)$, $w'$ y resultado tienen exactamente las mismas probabilidades, así que se
guardan juntas: \col{n} es cuántas observaciones hay en la celda y \col{n\_r1} cuántas
repararon de verdad. No se pierde nada. Para cualquier promedio, pondera por \col{n}; para
volver a una fila por observación, usa \texttt{tidyr::uncount(n)}, aunque crece mucho.

\begin{center}\small
\begin{adjustbox}{max width=\textwidth}\begin{tabular}{ll}
\toprule
Columna & Significado\\
\midrule
\col{rep}, \col{estimador} & réplica y con qué $\theta$ se calcularon las probabilidades\\
\col{regimen}, \col{tipo} & como siempre\\
\col{j}, \col{a}, \col{w} & el coche que se podía reparar (\col{w} es índice del grid)\\
\col{accidente} & 1 si el coche se descompuso (entonces \col{w\_next} $=-1$)\\
\col{w\_next} & índice del $w'$ con el que llegó al año siguiente\\
\col{o}, \col{salida} & resultado del año siguiente (códigos de la sección 2.4); \col{salida} $=1$ si se chatarrizó\\
\col{n}, \col{n\_r1} & observaciones en la celda y cuántas repararon de verdad\\
\col{prior}, \col{post\_d2}, \col{post\_d3} & las tres probabilidades de arriba\\
\bottomrule
\end{tabular}\end{adjustbox}
\end{center}

Ejemplo de la prueba de humo (grid chico, $w=-1.5+0.15\times$índice):
\begin{center}\footnotesize
\begin{adjustbox}{max width=\textwidth}\begin{tabular}{rlrrrrrrrrrrrrrr}
\toprule
rep & estimador & regimen & tipo & j & a & w & accidente & w\_next & o & salida & n & n\_r1 & prior & post\_d2 & post\_d3\\
\midrule
0 & verdad & 0 & 0 & 0 & 1 & 8 & 0 & 5 & 0 & 0 & 1 & 1 & 0.238 & 0.951 & 0.951\\
0 & verdad & 0 & 0 & 0 & 1 & 8 & 0 & 7 & 0 & 0 & 1 & 1 & 0.238 & 0.465 & 0.465\\
0 & verdad & 0 & 0 & 0 & 1 & 8 & 0 & 9 & 1 & 0 & 1 & 0 & 0.238 & 0.039 & 0.039\\
0 & verdad & 0 & 0 & 0 & 1 & 8 & 0 & 11 & 0 & 0 & 1 & 0 & 0.238 & 0.002 & 0.002\\
\bottomrule
\end{tabular}\end{adjustbox}
\end{center}
Se lee así: un light brown de 1 año con $w$ en el índice 8 tenía 24\% de probabilidad de ser
reparado. Si al año siguiente aparece con el desgaste más bajo (índice 5), casi seguro se
reparó (95\%). Si aparece con el desgaste más alto (11), casi seguro no se reparó (0.2\%).
Si su desgaste bajó un poco (7), el dato es ambiguo (47\%).

\paragraph{Tamaño.} Unas decenas de miles de celdas por réplica y estimador. Las 50 réplicas
son algunos millones de filas, en un archivo comprimido de algunos cientos de MB a lo más.
Léelo con \texttt{readr::read\_csv}, que descomprime solo.

% ===========================================================================
\section{Gráficas}

\subsection{Las que hice para el MC de octubre}

Están en \texttt{claude/niu/output/modelo\_tesis/reportes/mc\_base/} y las hace
\texttt{claude/niu/analisis/reporte\_mc.py}. Ese script apunta a las salidas viejas; lo copio
a \texttt{final/} cuando tengas el MC nuevo, si lo quieres usar de base.

\begin{center}\small
\begin{adjustbox}{max width=\textwidth}\begin{tabular}{p{3.4cm}p{3.6cm}p{7.6cm}}
\toprule
Figura & Base & Qué muestra\\
\midrule
fig1 sesgo & \texttt{parametros} + \texttt{gill\_parametros} & Sesgo relativo medio de los 20 parámetros comunes (21 en el MC de octubre, que tenía \col{tc\_sell\_inspect}), por estimador. Gillingham contra los tres diseños, con un zoom a los diseños e IC 95\% del sesgo medio.\\
fig2 reparación & \texttt{parametros} & Distribución entre réplicas de $\hat\kappa,\hat\sigma_\eta,\hat\sigma_r,\hat u_w,\hat\delta_j$ por diseño, con la verdad marcada.\\
fig3 cobertura & \texttt{parametros} & Cobertura del IC 95\% por parámetro y estimador, con la banda binomial de lo esperable con 50 réplicas.\\
fig4 errores estándar & \texttt{parametros} & sd entre réplicas contra el $se$ mediano: si caen en la diagonal, los errores estándar son correctos.\\
fig5 fallas & \texttt{resumen}, \texttt{fallas} & Qué (réplica, diseño) no tienen estimación.\\
fig6 precios por edad & \texttt{por\_edad\_mc} & Precio por edad y marca: verdad (línea) contra la media y la banda 5--95\% de las réplicas, por estimador.\\
fig7 reparación por edad & \texttt{por\_edad\_mc} & Pr(reparar) por edad; solo diseños, porque Gillingham no repara.\\
fig8 distribución & \texttt{por\_edad\_mc} & $q(j,a)$ por edad.\\
\bottomrule
\end{tabular}\end{adjustbox}
\end{center}

\subsection{Otras que podrían salir}

\paragraph{Del equilibrio verdadero (\texttt{teoria/}): el modelo.}
\begin{enumerate}
  \item \textbf{Respuesta de la reparación a su precio.} Pr(reparar) por edad, una línea por
  régimen (\texttt{por\_edad.csv}). Es la figura de la variable excluida: muestra que $R$
  mueve la reparación, y en qué edades más.
  \item \textbf{Desgaste medio por edad y régimen} (\col{w\_medio}). Muestra cómo la reparación
  barata mantiene los coches menos desgastados.
  \item \textbf{Modelo contra Gillingham.} Precio y $q$ por edad contra \col{P\_gill} y
  \col{q\_gill}. Responde si nuestra economía se parece a la danesa.
  \item \textbf{Precio de transacción contra precio de stock.} La diferencia indica qué coches
  se venden: si los que se venden están más desgastados que el promedio.
  \item \textbf{Agregados por régimen} (\texttt{resumen.csv}): reparación, accidentes,
  chatarreo y edad media contra $\zeta$.
\end{enumerate}

\paragraph{Del panel (\texttt{panel.csv.gz}): la identificación, vista en los datos.}
\begin{enumerate}
  \setcounter{enumi}{5}
  \item \textbf{La mezcla de Hu \& Xin.} Histograma de $w_{k+1}-w_{h,k}$ para coches que
  sobreviven, separado por $r=0$ y $r=1$. Salen dos normales centradas en $\delta$ y
  $\delta-\kappa$. Junto, el histograma sin separar, que es lo que ve el econometrista en los
  diseños 2 y 3. Es la mejor figura para explicar qué se identifica.
  \item \textbf{La misma mezcla en los tres regímenes.} El peso de cada normal cambia con $R$
  y las dos densidades no se mueven. Es el Supuesto 7 en una figura.
  \item \textbf{Reparación contra desgaste:} tasa de $r=1$ por valor de $w$ (y edad). Muestra
  quién repara.
  \item \textbf{Accidentes por edad y por $w$:} tasa observada contra la curva
  $s(j,a,w)$.
  \item \textbf{Distribución de $w$ por edad} (violines o percentiles): cómo se dispersa el
  desgaste con los años.
\end{enumerate}

\paragraph{Del Monte Carlo: los resultados.}
\begin{enumerate}
  \setcounter{enumi}{10}
  \item \textbf{El costo de no observar $r$.} Para cada parámetro, el cociente
  $sd(\text{diseño }2)/sd(\text{diseño }1)$ y $sd(\text{diseño }3)/sd(\text{diseño }1)$.
  Es la pregunta 2 de la tesis en una sola figura (barras por parámetro).
  \item \textbf{Correlación entre $\hat\kappa$ y $\hat\sigma_\eta$} (o $\hat\kappa$ y
  $\hat\delta_j$) a través de las réplicas: dispersión por diseño. Sin $r$, la mezcla hace que
  se compensen; la elipse se alarga.
  \item \textbf{Tasa de reparación estimada contra verdadera} por réplica
  (\col{reparacion\_est} contra \col{reparacion\_verdad} en \texttt{resumen\_reps}).
  \item \textbf{Error en precios por diseño} (\col{P\_rmse}): no observar $r$ ni el motivo,
  ¿cuánto empeora los precios implicados?
  \item \textbf{Lo que Gillingham ``cree'' del mercado:} agregados implicados por su
  $\hat\theta$ (\texttt{gill\_resumen}, columnas \col{*\_est}) contra los verdaderos
  (\texttt{teoria/resumen.csv}). Por ejemplo: atribuye a chatarreo voluntario lo que en
  realidad es desgaste.
  \item \textbf{Disposición a pagar en DKK:} $u_0/\mu$ y $T_b/\mu$ estimados contra los
  verdaderos. Son más interpretables que los parámetros en utils, y el sesgo de Gillingham en
  $\mu$ se nota más.
\end{enumerate}

\paragraph{De la inferencia de $r$ (\texttt{inferencia\_reparacion}): qué tanto se infiere la
reparación de cada coche.}
\begin{enumerate}
  \setcounter{enumi}{16}
  \item \textbf{Curvas ROC.} Para cada estimador y nivel de información (\col{prior},
  \col{post\_d2}, \col{post\_d3}), la tasa de verdaderos positivos contra la de falsos
  positivos al mover el umbral. Los positivos de una celda son \col{n\_r1} y los negativos
  \col{n}$-$\col{n\_r1}. El área bajo la curva (AUC) resume en un número cuánto se infiere $r$.
  Comparar \col{prior} con \col{post\_d2} muestra cuánto informa ver $w'$.
  \item \textbf{Calibración.} Agrupa por deciles de la probabilidad y grafica la probabilidad
  media contra la tasa real \col{n\_r1}/\col{n}. Con \col{verdad} cae en la diagonal; con
  $\hat\theta$, la distancia a la diagonal es el error de estimación.
  \item \textbf{Costo de la estimación.} AUC (o Brier) con \col{verdad} contra con
  \col{diseno\_2} y \col{diseno\_3}, réplica por réplica: ¿se pierde capacidad de inferir $r$
  por no conocer $\theta$?
  \item \textbf{Qué observaciones son ambiguas.} \col{post\_d2} contra \col{w\_next}$-$\col{w},
  coloreado por el $r$ verdadero: los coches cuyo desgaste cambió cerca del punto medio entre
  $\delta$ y $\delta-\kappa$ son los que no se pueden clasificar.
  \item \textbf{Lo que se pierde en las salidas.} Entre las celdas con \col{salida = 1},
  \col{post\_d3} contra \col{post\_d2}: cuánto se pierde por no ver el $w'$ de los coches
  chatarrizados (el costo del diseño 3 frente al 2, observación por observación).
  \item \textbf{Inferencia por edad y régimen.} AUC o Brier por \col{a} y por \col{regimen}.
\end{enumerate}

\end{document}
